% Compile with: uplatex file.tex && dvipdfmx file.dvi
%           or: latexmk -uplatex file.tex
\documentclass[dvipdfmx]{article}
\usepackage{amsmath, amssymb}
\usepackage{tikz}

\title{scenario\_line\_balancing}
\date{\today}

\begin{document}
\maketitle

From the total work-element time $T$ and the cycle time $c$ (the time allowed per station), define the minimum number of stations $n=\lceil T/c\rceil$ (ceiling function, since fewer than $T/c$ stations cannot suffice).

$n = ceiling(T / c)$

\begin{equation}
  n = ceiling(T / c)
  \label{eq:E1}
\end{equation}

Define the line efficiency $\eta=T/(nc)$ (the ratio of the ideal total working time $T$ to the actual total available time $nc$).

$\eta = T / (n c)$

\begin{equation}
  \eta = T / (n c)
  \label{eq:E2}
\end{equation}

As a concrete example, substitute the total work time $T=88$ into the station-count expression.

In E1, substitute 88 for T

\begin{equation}
  n = \left \lceil \frac{88}{c} \right \rceil
  \label{eq:E3}
\end{equation}

Substitute the cycle time $c=15$ as well, to fix the minimum number of stations $n$ as a number.

In E3, substitute 15 for c

\begin{equation}
  n = 6
  \label{eq:E4}
\end{equation}

Substitute the same total work time $T=88$, this time into the efficiency expression.

In E2, substitute 88 for T

\begin{equation}
  \eta = \frac{88}{\left \lceil \frac{88}{c} \right \rceil c}
  \label{eq:E5}
\end{equation}

Substitute the same cycle time $c=15$, to fix the line efficiency $\eta$ as a number.

In E5, substitute 15 for c

\begin{equation}
  \eta = \frac{44}{45}
  \label{eq:E6}
\end{equation}

Compute the balance delay (idle ratio) $d=1-\eta$, plugging in the value of $\eta$ fixed in the previous step directly.

$d = 1 - E6$

\begin{equation}
  d = \frac{1}{45}
  \label{eq:E7}
\end{equation}

\end{document}